\documentclass[8pt]{beamer}
\usetheme{Madrid}
\usecolortheme{default}
\usepackage{amsmath, amssymb, graphicx, array, multirow, booktabs, xcolor, tikz}
\usetikzlibrary{shapes,arrows,positioning,fit,backgrounds}

% Compact font settings
\setbeamerfont{title}{size=\Large}
\setbeamerfont{subtitle}{size=\small}
\setbeamerfont{author}{size=\scriptsize}
\setbeamerfont{date}{size=\scriptsize}
\setbeamerfont{frametitle}{size=\large}
\setbeamerfont{framesubtitle}{size=\normalsize}
\setbeamerfont{enumerate item}{size=\footnotesize}
\setbeamerfont{itemize item}{size=\footnotesize}
\setbeamerfont{block title}{size=\footnotesize}

% Compact spacing
\setlength{\itemsep}{0.08ex}
\setlength{\parskip}{0.08ex}
\setlength{\leftmargini}{0.3cm}
\setbeamersize{text margin left=3mm, text margin right=3mm}
\addtobeamertemplate{frame begin}{\vspace{-0.4cm}}{}

% Colors
\definecolor{myblue}{RGB}{0,0,255}
\definecolor{myred}{RGB}{255,0,0}
\definecolor{mygreen}{RGB}{0,128,0}
\definecolor{myorange}{RGB}{255,165,0}
\definecolor{mypurple}{RGB}{128,0,128}
\definecolor{mygray}{RGB}{128,128,128}
\definecolor{mygold}{RGB}{255,215,0}

\title{Responsible \& Ethical AI: 50 Advanced Questions}
\subtitle{Deep-dive questions: theory, technical mechanisms, and regulation}
\author{GARP RAI Exam Preparation}
\date{}

\begin{document}

\begin{frame}
	\titlepage
\end{frame}

% ===== FRAMEWORKS =====

\begin{frame}[shrink=6]{Q1: Rule Consequentialism Collapse - Advanced}
\fontsize{6.5}{7.5}\selectfont
\textbf{QUESTION: A rule consequentialist refuses to kill one patient to save five because a rule against killing maximizes long-run utility (it preserves trust in medicine). Critics say the theory ``collapses into deontology.'' What is the core of that objection?}

\vspace{0.1cm}
\textbf{OPTIONS:}
\begin{enumerate}[A.]
\item Rule consequentialists never consider outcomes, so they were deontologists from the start
\item If the rule must be obeyed even where breaking it would yield more utility, rule-adherence works as an independent constraint, which is exactly what deontologists claim
\item Deontologists also choose their rules solely by maximizing aggregate utility, so the two share one justification
\item Rules only matter case by case, so the theory reduces to act utilitarianism
\end{enumerate}

\vspace{0.1cm}
\textcolor{myblue}{\textbf{ANSWER: B}}

\vspace{0.1cm}
\textbf{DETAILED EXPLANATION:}
\begin{itemize}
\item \textbf{Why B is correct:} Rules are justified by their tendency to maximize utility, but acts are judged by conformity to the rule, not by case-specific outcomes. When conformity overrides a better isolated outcome, the rule behaves like a deontological red line.
\item \textbf{Why A is wrong:} Rules are selected by their consequences, so outcomes still matter at the rule level
\item \textbf{Why C is wrong:} Deontologists ground rules in duties and rights, not in utility maximization
\item \textbf{Why D is wrong:} That describes act utilitarianism, the view rule consequentialism was designed to repair
\end{itemize}

\textcolor{myred}{\textbf{EXAM TRAP:}} Rule consequentialism: consequences pick the RULES, rules (not acts) judge conduct, so critics say it ends up deontological!
\end{frame}

\begin{frame}[shrink=6]{Q2: Kant vs Contemporary Deontology - Advanced}
\fontsize{6.5}{7.5}\selectfont
\textbf{QUESTION: An AI system could deceive users and produce large net benefits. Which analysis correctly contrasts Kant with most contemporary deontologists?}

\vspace{0.1cm}
\textbf{OPTIONS:}
\begin{enumerate}[A.]
\item Kant permits lying when benefits are large; contemporary deontologists forbid it absolutely
\item Both ignore consequences entirely, which is what separates them from virtue ethicists
\item Contemporary deontologists are act utilitarians who use rights only as tie-breakers
\item Kant treats lying as wrong irrespective of consequences; most contemporary deontologists give consequences weight but hold that some red lines cannot be crossed even for good overall outcomes
\end{enumerate}

\vspace{0.1cm}
\textcolor{myblue}{\textbf{ANSWER: D}}

\vspace{0.1cm}
\textbf{DETAILED EXPLANATION:}
\begin{itemize}
\item \textbf{Why D is correct:} Kant places no moral weight on consequences. Contemporary deontologists care about consequences, but rights, rules, and principles act as constraints that go beyond them.
\item \textbf{Why A is wrong:} This reverses the positions; Kant is the stricter view
\item \textbf{Why B is wrong:} Most contemporary deontologists do grant consequences moral weight
\item \textbf{Why C is wrong:} Consequences are one consideration for them, not the only one, so they are not utilitarians
\end{itemize}

\textcolor{myred}{\textbf{EXAM TRAP:}} Kant = consequences carry no weight; modern deontologists = consequences count, but RED LINES remain!
\end{frame}

\begin{frame}[shrink=6]{Q3: Universalizability Test - Advanced}
\fontsize{6.5}{7.5}\selectfont
\textbf{QUESTION: A firm adopts the maxim ``we may mislead customers whenever it is commercially convenient.'' Which reasoning applies the categorical imperative most accurately?}

\vspace{0.1cm}
\textbf{OPTIONS:}
\begin{enumerate}[A.]
\item It is rejected because average outcomes across firms would be worse in a utility calculation
\item It is rejected because a virtuous person would feel ashamed of it
\item It is acceptable as long as competitors do not adopt the same maxim
\item The maxim cannot be willed as a universal law without contradiction: if everyone misled whenever convenient, trust in statements would collapse and the strategy would defeat itself
\end{enumerate}

\vspace{0.1cm}
\textcolor{myblue}{\textbf{ANSWER: D}}

\vspace{0.1cm}
\textbf{DETAILED EXPLANATION:}
\begin{itemize}
\item \textbf{Why D is correct:} The test asks whether a maxim can be universalized without contradiction, regardless of consequences. Deception parasitically depends on a background practice of truth-telling that universal deception would destroy.
\item \textbf{Why A is wrong:} Cost-benefit aggregation is consequentialist, not the Kantian test
\item \textbf{Why B is wrong:} Asking what a virtuous agent would do is virtue ethics
\item \textbf{Why C is wrong:} Exempting oneself from a rule that others follow is exactly what universalizability forbids
\end{itemize}

\textcolor{myred}{\textbf{EXAM TRAP:}} Categorical imperative = could the maxim be willed as universal law WITHOUT CONTRADICTION?
\end{frame}

\begin{frame}[shrink=6]{Q4: Undercutting the Utilitarian Metric - Advanced}
\fontsize{6.5}{7.5}\selectfont
\textbf{QUESTION: Utilitarianism is praised for providing a single, quantifiable, impartial metric. Which statement best explains why this advantage is undercut in practice?}

\vspace{0.1cm}
\textbf{OPTIONS:}
\begin{enumerate}[A.]
\item Utilitarianism is backward-looking, so it cannot evaluate future AI deployments
\item Utilitarianism requires partiality toward those closest to the decision maker
\item What counts as utility (pleasure, happiness, desire satisfaction) is subjective and hard to quantify, and consequences are unpredictable with no clear stopping point for tracing consequences of consequences
\item Utility can only be defined as pleasure, which cannot be measured
\end{enumerate}

\vspace{0.1cm}
\textcolor{myblue}{\textbf{ANSWER: C}}

\vspace{0.1cm}
\textbf{DETAILED EXPLANATION:}
\begin{itemize}
\item \textbf{Why C is correct:} The metric is single in form but contestable in content, and the scope of the calculation is open-ended, so apparent objectivity is fragile.
\item \textbf{Why A is wrong:} Utilitarianism is forward-looking and focused on end consequences
\item \textbf{Why B is wrong:} It aims to be impartial and objective, not partial
\item \textbf{Why D is wrong:} Utility may also be defined as happiness or satisfaction of desires
\end{itemize}

\textcolor{myred}{\textbf{EXAM TRAP:}} Utilitarianism's weakness: choice of metric is subjective AND where to stop the calculation is unclear!
\end{frame}

\begin{frame}[shrink=6]{Q5: Wrongs Without Harm - Advanced}
\fontsize{6.5}{7.5}\selectfont
\textbf{QUESTION: The module says violating privacy is wrong even when it causes no harm, comparing it to cheating on a spouse who never finds out. Which foundation best supports that claim?}

\vspace{0.1cm}
\textbf{OPTIONS:}
\begin{enumerate}[A.]
\item Privacy violations can treat persons merely as means to the violator's ends and infringe moral and legal rights, so they wrong someone independently of tangible harm
\item They are wrong only in expectation, because they raise the probability of later harm
\item They are wrong only where a statute recognizes a right to privacy
\item They are wrong because they reduce aggregate welfare
\end{enumerate}

\vspace{0.1cm}
\textcolor{myblue}{\textbf{ANSWER: A}}

\vspace{0.1cm}
\textbf{DETAILED EXPLANATION:}
\begin{itemize}
\item \textbf{Why A is correct:} A wrong is a violation of rights or of a person's standing; harm is a separate notion. Treating people instrumentally, as things and not autonomous beings, is wrong even if nothing bad follows.
\item \textbf{Why B is wrong:} This makes the wrong depend on harm, which the module explicitly denies
\item \textbf{Why C is wrong:} The right to privacy is a moral right whether or not a law or era recognizes it
\item \textbf{Why D is wrong:} Aggregate welfare is consequentialist reasoning and again depends on harm
\end{itemize}

\textcolor{myred}{\textbf{EXAM TRAP:}} WRONGS can exist without HARM: rights violations and treating people as mere means!
\end{frame}

\begin{frame}[shrink=6]{Q6: Virtue Ethics Guidance Problem - Advanced}
\fontsize{6.5}{7.5}\selectfont
\textbf{QUESTION: Critics say virtue ethics gives less guidance than duty-based or consequentialist approaches because virtues can conflict. What reply does the module give?}

\vspace{0.1cm}
\textbf{OPTIONS:}
\begin{enumerate}[A.]
\item Virtues never actually conflict once each is properly understood
\item Virtue ethics supplies a fixed ranking of virtues that resolves conflicts mechanically
\item Guidance is unnecessary because character is fixed in early childhood
\item Practical wisdom helps navigate hard cases, and rival theories face parallel problems: duties can conflict and consequences are not always comparable
\end{enumerate}

\vspace{0.1cm}
\textcolor{myblue}{\textbf{ANSWER: D}}

\vspace{0.1cm}
\textbf{DETAILED EXPLANATION:}
\begin{itemize}
\item \textbf{Why D is correct:} Virtue ethicists claim they are no worse off: every framework must weigh competing considerations, and practical wisdom is the faculty for doing so in context.
\item \textbf{Why A is wrong:} The module accepts that virtues can conflict; its point is that all theories face conflict
\item \textbf{Why B is wrong:} There is no mechanical ranking; practical wisdom is exercised case by case
\item \textbf{Why C is wrong:} Virtues are built through practice, habit, and exemplars over a lifetime
\end{itemize}

\textcolor{myred}{\textbf{EXAM TRAP:}} Virtue ethics reply = PRACTICAL WISDOM, plus 'you have conflicts too' (duties clash, consequences incomparable)!
\end{frame}

\begin{frame}[shrink=6]{Q7: Virtue Ethics and Building AI - Advanced}
\fontsize{6.5}{7.5}\selectfont
\textbf{QUESTION: Some scholars argue that to build ethical AI we must build it ``in a virtue ethics way.'' What is the argument?}

\vspace{0.1cm}
\textbf{OPTIONS:}
\begin{enumerate}[A.]
\item Morality is so complex that it may never be coded top-down, so systems should learn from experience and habituation, the way children acquire morality
\item Virtue ethics provides an explicit rule set that can be encoded as hard constraints
\item Virtue ethics equates goodness with total utility, which a loss function can optimize
\item AI systems have emotions and appetites that must be moderated
\end{enumerate}

\vspace{0.1cm}
\textcolor{myblue}{\textbf{ANSWER: A}}

\vspace{0.1cm}
\textbf{DETAILED EXPLANATION:}
\begin{itemize}
\item \textbf{Why A is correct:} The virtue approach is bottom-up: character is formed through practice, habit, and modeling exemplars, which parallels learning from experience.
\item \textbf{Why B is wrong:} Rule sets are the deontological, top-down route the argument says may fail
\item \textbf{Why C is wrong:} Utility maximization is consequentialist
\item \textbf{Why D is wrong:} The point is about how morality is learned, not about machine emotions
\end{itemize}

\textcolor{myred}{\textbf{EXAM TRAP:}} Virtue-ethics AI = learn from experience and habit (BOTTOM-UP), because top-down coding of morality may be impossible!
\end{frame}

\begin{frame}[shrink=6]{Q8: Three Lenses on Privacy Decisions - Advanced}
\fontsize{6.5}{7.5}\selectfont
\textbf{QUESTION: Which mapping of the three frameworks onto decisions about collecting and keeping personal data matches the module?}

\vspace{0.1cm}
\textbf{OPTIONS:}
\begin{enumerate}[A.]
\item Consequentialist: the right to privacy. Deontological: worst-case scenarios. Virtue: data minimization
\item Consequentialist: what would a responsible company do. Deontological: best/worst cases. Virtue: data minimization
\item All three reduce to compliance with the applicable law
\item Consequentialist: best/worst cases, data sensitivity, confidence in keeping it safe. Deontological: the right to privacy and principles like data minimization. Virtue: what would a responsible company do
\end{enumerate}

\vspace{0.1cm}
\textcolor{myblue}{\textbf{ANSWER: D}}

\vspace{0.1cm}
\textbf{DETAILED EXPLANATION:}
\begin{itemize}
\item \textbf{Why D is correct:} Consequentialism asks about outcomes and risk, deontology about rights and rules, and virtue ethics about the character of a responsible organization.
\item \textbf{Why A is wrong:} Rights and principles are deontological, so the roles are swapped
\item \textbf{Why B is wrong:} Character questions belong to virtue ethics, so the roles are scrambled
\item \textbf{Why C is wrong:} The module says ethics goes beyond minimal legal requirements
\end{itemize}

\textcolor{myred}{\textbf{EXAM TRAP:}} Consequences = risk and sensitivity; Deontology = rights and minimization; Virtue = what a responsible company would do!
\end{frame}

\begin{frame}[shrink=6]{Q9: Ethics and Law - Advanced}
\fontsize{6.5}{7.5}\selectfont
\textbf{QUESTION: The module calls ethics ``more ambitious than the law'' and says it helps distinguish just from unjust laws. Which inference is consistent with this?}

\vspace{0.1cm}
\textbf{OPTIONS:}
\begin{enumerate}[A.]
\item Whatever is legal is ethically permissible, because law fully encodes social morality
\item Ethics becomes redundant once laws like the GDPR exist
\item Abandoning a friend in need is legal yet immoral; because ethics offers standards independent of statutes, it can ground, inform, and criticize law
\item Law sets aspirational ideals of the good life while ethics sets minimum requirements
\end{enumerate}

\vspace{0.1cm}
\textcolor{myblue}{\textbf{ANSWER: C}}

\vspace{0.1cm}
\textbf{DETAILED EXPLANATION:}
\begin{itemize}
\item \textbf{Why C is correct:} Laws typically set minimal behavioral requirements so institutions function. Ethics asks what kind of society and life we want and therefore can shape and evaluate laws.
\item \textbf{Why A is wrong:} Legal does not imply ethical; the friend example shows the gap
\item \textbf{Why B is wrong:} The module says laws are also needed but do not replace ethics
\item \textbf{Why D is wrong:} This reverses the roles of law (minimal) and ethics (ambitious)
\end{itemize}

\textcolor{myred}{\textbf{EXAM TRAP:}} Law = MINIMUM requirements; ethics = more ambitious, and can judge whether laws are just!
\end{frame}

% ===== MEDICAL ETHICS =====

\begin{frame}[shrink=6]{Q10: Moral vs Medical Questions - Advanced}
\fontsize{6.5}{7.5}\selectfont
\textbf{QUESTION: Why does the module use the ventilator and organ procurement from heart-beating bodies with non-functioning brains as a case for formal ethics structures?}

\vspace{0.1cm}
\textbf{OPTIONS:}
\begin{enumerate}[A.]
\item Clinicians lacked the technical skill to operate ventilators safely
\item The case proved that consequentialism is the only workable framework in medicine
\item Whether to take such organs is a moral, not medical, question, so resolving it should not rest solely on clinicians, whose expertise is maintaining health
\item Regulators required hospitals to adopt AI ethics codes
\end{enumerate}

\vspace{0.1cm}
\textcolor{myblue}{\textbf{ANSWER: C}}

\vspace{0.1cm}
\textbf{DETAILED EXPLANATION:}
\begin{itemize}
\item \textbf{Why C is correct:} New technology created dilemmas that medical training does not resolve. This drove ethics frameworks, committees, and codes, and is offered as a parallel for engineers and executives facing AI dilemmas.
\item \textbf{Why A is wrong:} The issue was moral, not a shortfall of technical skill
\item \textbf{Why B is wrong:} The module says frameworks are plural and complementary
\item \textbf{Why D is wrong:} AI ethics codes are not part of the medical case
\end{itemize}

\textcolor{myred}{\textbf{EXAM TRAP:}} Ventilator lesson: moral questions are not medical questions, so experts should not decide them alone!
\end{frame}

% ===== PRINCIPLES =====

\begin{frame}[shrink=6]{Q11: Nonmaleficence and Risk Allocation - Advanced}
\fontsize{6.5}{7.5}\selectfont
\textbf{QUESTION: Which scenario is MOST defensible under the module's account of nonmaleficence, which considers who decides and who bears the brunt of harm?}

\vspace{0.1cm}
\textbf{OPTIONS:}
\begin{enumerate}[A.]
\item Very risky clinical research on patients with a serious ailment who stand to benefit if it succeeds
\item The same very risky research on healthy volunteers
\item The same research on patients whose ailment the research cannot cure
\item An algorithm whose risks fall on people who gain nothing, imposed by people facing no risk
\end{enumerate}

\vspace{0.1cm}
\textcolor{myblue}{\textbf{ANSWER: A}}

\vspace{0.1cm}
\textbf{DETAILED EXPLANATION:}
\begin{itemize}
\item \textbf{Why A is correct:} Risk is more acceptable when imposed on people who stand to benefit. Nonmaleficence does not forbid all harm; it forbids unjustified risk that outweighs benefits.
\item \textbf{Why B is wrong:} Healthy subjects bear risk without prospect of benefit
\item \textbf{Why C is wrong:} No possible benefit to those bearing the risk
\item \textbf{Why D is wrong:} This is the module's own analogy for an unacceptable algorithmic risk
\end{itemize}

\textcolor{myred}{\textbf{EXAM TRAP:}} Nonmaleficence is not zero-risk: ask WHO benefits and WHO bears the risk!
\end{frame}

\begin{frame}[shrink=6]{Q12: Unjustified Risk Imposition - Advanced}
\fontsize{6.5}{7.5}\selectfont
\textbf{QUESTION: A firm deploys an untested algorithm that could harm users. No harm has yet occurred. Under nonmaleficence:}

\vspace{0.1cm}
\textbf{OPTIONS:}
\begin{enumerate}[A.]
\item The firm may already be acting wrongly, because imposing unjustified or unnecessary risk on others is itself covered by the obligation
\item No violation exists until measurable harm materializes
\item Deployment is acceptable provided expected profit exceeds expected harm to the firm
\item The principle applies only to intentional injury
\end{enumerate}

\vspace{0.1cm}
\textcolor{myblue}{\textbf{ANSWER: A}}

\vspace{0.1cm}
\textbf{DETAILED EXPLANATION:}
\begin{itemize}
\item \textbf{Why A is correct:} The module says nonmaleficence includes a prohibition on imposing risks of harm that are not justified or that outweigh benefits, for example subjecting people to untested harmful algorithms.
\item \textbf{Why B is wrong:} The obligation reaches risk imposition, not only realized harm
\item \textbf{Why C is wrong:} Justification must weigh risks to those exposed, not only to the firm
\item \textbf{Why D is wrong:} The module includes unnecessary and unjustified risk, not only intentional injury
\end{itemize}

\textcolor{myred}{\textbf{EXAM TRAP:}} Nonmaleficence covers RISK IMPOSITION, not only realized harm!
\end{frame}

\begin{frame}[shrink=6]{Q13: Beneficence Is Not Merely Consequentialist - Advanced}
\fontsize{6.5}{7.5}\selectfont
\textbf{QUESTION: Which statement about beneficence is correct according to the module?}

\vspace{0.1cm}
\textbf{OPTIONS:}
\begin{enumerate}[A.]
\item It is purely a consequentialist idea with no place in duty-based ethics
\item It is unlimited: each agent must do everything possible to relieve suffering
\item Deontological frameworks also include beneficence as a duty above what utility maximization requires, but it is limited by scarce resources, competing obligations, reasonableness, and demandingness
\item It is equivalent to nonmaleficence stated positively
\end{enumerate}

\vspace{0.1cm}
\textcolor{myblue}{\textbf{ANSWER: C}}

\vspace{0.1cm}
\textbf{DETAILED EXPLANATION:}
\begin{itemize}
\item \textbf{Why C is correct:} Beneficence requires positive steps to help others, a key deontological duty, but ethics cannot demand unlimited sacrifice.
\item \textbf{Why A is wrong:} This is the common misunderstanding the module corrects
\item \textbf{Why B is wrong:} No individual can alleviate all suffering; reasonable constraints apply
\item \textbf{Why D is wrong:} Nonmaleficence forbids injuring; beneficence commands positive help
\end{itemize}

\textcolor{myred}{\textbf{EXAM TRAP:}} Beneficence is ALSO a deontological duty, bounded by demandingness and reasonableness!
\end{frame}

\begin{frame}[shrink=6]{Q14: Beneficence vs Autonomy - Advanced}
\fontsize{6.5}{7.5}\selectfont
\textbf{QUESTION: A wellness app decides to override a competent adult's declared choice because it computes that another option would benefit them. What does the module imply?}

\vspace{0.1cm}
\textbf{OPTIONS:}
\begin{enumerate}[A.]
\item Autonomy usually trumps unwanted benefits, with exceptions such as public health worries or doubts about the person's decision-making capacity
\item Beneficence always outranks autonomy when the benefit is large
\item Autonomy is absolute and can never be overridden, even for children
\item Surrogate decision makers are needed for every adult decision
\end{enumerate}

\vspace{0.1cm}
\textcolor{myblue}{\textbf{ANSWER: A}}

\vspace{0.1cm}
\textbf{DETAILED EXPLANATION:}
\begin{itemize}
\item \textbf{Why A is correct:} People have a right to decide what is best for them. Only public-health concerns or worries about capacity typically justify overriding, and then a surrogate acts in best interests.
\item \textbf{Why B is wrong:} Unwanted benefits that disrespect agency are precluded by the module
\item \textbf{Why C is wrong:} The module allows surrogates for those lacking capacity (children, unconscious patients)
\item \textbf{Why D is wrong:} Surrogates are for people lacking capacity, not for every adult
\end{itemize}

\textcolor{myred}{\textbf{EXAM TRAP:}} Autonomy usually TRUMPS unwanted beneficence; exceptions: public health and lack of capacity!
\end{frame}

\begin{frame}[shrink=6]{Q15: Autonomy and Manipulative Design - Advanced}
\fontsize{6.5}{7.5}\selectfont
\textbf{QUESTION: Which system MOST clearly contravenes autonomy as defined in the module (informed, un-coerced choice; technology should advance the user's goals, not third parties')?}

\vspace{0.1cm}
\textbf{OPTIONS:}
\begin{enumerate}[A.]
\item A decision aid that shows its recommendation and lets users freely override it
\item A recommender that exploits psychological weaknesses to steer purchases benefiting the platform, even though users technically click accept
\item A consent flow that discloses data uses and allows withdrawal at any time
\item A guardian-facing tool that helps a legal guardian act in a child's best interests
\end{enumerate}

\vspace{0.1cm}
\textcolor{myblue}{\textbf{ANSWER: B}}

\vspace{0.1cm}
\textbf{DETAILED EXPLANATION:}
\begin{itemize}
\item \textbf{Why B is correct:} Coercion, deception, manipulation, and undue influence all undermine autonomy. Nominal clicking does not make manipulated choices autonomous.
\item \textbf{Why A is wrong:} Overridable recommendations preserve choice
\item \textbf{Why C is wrong:} Meaningful consent and withdrawal support autonomy
\item \textbf{Why D is wrong:} Surrogate decision-making for those lacking capacity is a recognized exception
\end{itemize}

\textcolor{myred}{\textbf{EXAM TRAP:}} Autonomy is undermined by manipulation even when the user 'agrees'!
\end{frame}

\begin{frame}[shrink=6]{Q16: Conceptions of Justice - Advanced}
\fontsize{6.5}{7.5}\selectfont
\textbf{QUESTION: A firm remedies harm from a biased hiring tool by meeting rejected applicants, acknowledging the harm, and compensating them. Which conception of justice is this?}

\vspace{0.1cm}
\textbf{OPTIONS:}
\begin{enumerate}[A.]
\item Procedural justice: fair, impartial processes
\item Restorative justice: repairing harms by reconciling victims and offenders
\item Distributive justice: equitable allocation of benefits and burdens
\item Interactional justice: respect and fairness in dealings between individuals, with no repair element
\end{enumerate}

\vspace{0.1cm}
\textcolor{myblue}{\textbf{ANSWER: B}}

\vspace{0.1cm}
\textbf{DETAILED EXPLANATION:}
\begin{itemize}
\item \textbf{Why B is correct:} Restorative justice aims to repair harm. The scenario centers on acknowledgment and remedy after harm has occurred.
\item \textbf{Why A is wrong:} Procedural concerns the fairness of the decision process itself
\item \textbf{Why C is wrong:} Distributive concerns allocation, not repair
\item \textbf{Why D is wrong:} Interactional concerns respectful treatment; the key here is repair
\end{itemize}

\textcolor{myred}{\textbf{EXAM TRAP:}} Restorative = REPAIR harm; procedural = fair process; distributive = allocation; interactional = respect; social = just institutions!
\end{frame}

\begin{frame}[shrink=6]{Q17: Competing Principles of Justice - Advanced}
\fontsize{6.5}{7.5}\selectfont
\textbf{QUESTION: One triage algorithm allocates scarce resources by need, another by merit. Both claim to be just. Which statement is most accurate?}

\vspace{0.1cm}
\textbf{OPTIONS:}
\begin{enumerate}[A.]
\item They reflect competing principles of justice (need, meritocracy, egalitarianism, utilitarianism) that can conflict, so legitimacy depends on decisions being justifiable to all, not on proving one principle correct
\item Only need-based allocation counts as justice; merit-based allocation is discrimination
\item Justice is purely procedural, so allocation principles are irrelevant
\item Broad agreement exists on which goods justice must distribute
\end{enumerate}

\vspace{0.1cm}
\textcolor{myblue}{\textbf{ANSWER: A}}

\vspace{0.1cm}
\textbf{DETAILED EXPLANATION:}
\begin{itemize}
\item \textbf{Why A is correct:} The module notes disagreement about what a just distribution is and which goods justice concerns. It notes broad agreement on discrimination for morally irrelevant traits.
\item \textbf{Why B is wrong:} Merit is a recognized justice principle; it is not automatically discrimination
\item \textbf{Why C is wrong:} Distributive justice concerns allocation and is a distinct conception
\item \textbf{Why D is wrong:} The module says people disagree on the goods (resources, opportunities, power, welfare)
\end{itemize}

\textcolor{myred}{\textbf{EXAM TRAP:}} Justice principles COMPETE; legitimacy requires decisions justifiable to all!
\end{frame}

\begin{frame}[shrink=6]{Q18: Local and Model-Agnostic Explainability - Advanced}
\fontsize{6.5}{7.5}\selectfont
\textbf{QUESTION: A regulator wants to know why one applicant was denied by a black-box ensemble the provider will not restructure. Which type of explanation fits best?}

\vspace{0.1cm}
\textbf{OPTIONS:}
\begin{enumerate}[A.]
\item A global explanation using rules tailored to decision-tree architecture
\item A local, model-agnostic explanation of that specific decision
\item A global, model-agnostic summary of behavior across all inputs
\item No explanation can exist for a black-box model
\end{enumerate}

\vspace{0.1cm}
\textcolor{myblue}{\textbf{ANSWER: B}}

\vspace{0.1cm}
\textbf{DETAILED EXPLANATION:}
\begin{itemize}
\item \textbf{Why B is correct:} Local explainability targets a single prediction. Model-agnostic methods work with any architecture, while model-specific techniques depend on the model type.
\item \textbf{Why A is wrong:} Global, and tree-specific rules do not apply to an ensemble
\item \textbf{Why C is wrong:} Global summaries do not explain one case
\item \textbf{Why D is wrong:} Explanation methods exist, and auditing tools are being developed for black boxes
\end{itemize}

\textcolor{myred}{\textbf{EXAM TRAP:}} Single decision = LOCAL; any architecture = MODEL-AGNOSTIC!
\end{frame}

\begin{frame}[shrink=6]{Q19: Counterfactual Explanations and Audience - Advanced}
\fontsize{6.5}{7.5}\selectfont
\textbf{QUESTION: Which statement best captures what counterfactual explanations offer and the caveat about explanations generally?}

\vspace{0.1cm}
\textbf{OPTIONS:}
\begin{enumerate}[A.]
\item They describe overall behavior across all inputs; one explanation fits every audience
\item They prove the model is unbiased; the audience is irrelevant
\item They state conditions under which the outcome would have differed (for example, \$10,000 more savings), giving actionable insight; what counts as an explanation depends on to whom it is owed
\item They disclose the training dataset; only experts need explanations
\end{enumerate}

\vspace{0.1cm}
\textcolor{myblue}{\textbf{ANSWER: C}}

\vspace{0.1cm}
\textbf{DETAILED EXPLANATION:}
\begin{itemize}
\item \textbf{Why C is correct:} Counterfactuals contrast the actual decision with a hypothetical one. The module stresses explanations for experts differ from those for regulators or ordinary citizens.
\item \textbf{Why A is wrong:} That describes global explainability
\item \textbf{Why B is wrong:} They illuminate variable influence, not the absence of bias
\item \textbf{Why D is wrong:} They concern decision conditions, not disclosure of training data
\end{itemize}

\textcolor{myred}{\textbf{EXAM TRAP:}} Counterfactual = 'what would have changed the decision?'; explanations depend on the AUDIENCE!
\end{frame}

% ===== BIAS AND FAIRNESS =====

\begin{frame}[shrink=6]{Q20: Bias Is Not Unfairness - Advanced}
\fontsize{6.5}{7.5}\selectfont
\textbf{QUESTION: Which statement is consistent with the module's view of algorithmic bias?}

\vspace{0.1cm}
\textbf{OPTIONS:}
\begin{enumerate}[A.]
\item Some biases are justifiable, and a statistically or legally unbiased algorithm can still be unethical (for example, one that charges exorbitant fees to everyone)
\item Any deviation from parity is unethical, so an unbiased algorithm is always ethical
\item A well-designed AI should be entirely objective and value-free
\item Bias arises only from prejudiced designers
\end{enumerate}

\vspace{0.1cm}
\textcolor{myblue}{\textbf{ANSWER: A}}

\vspace{0.1cm}
\textbf{DETAILED EXPLANATION:}
\begin{itemize}
\item \textbf{Why A is correct:} Bias is deviation from a desired norm or intended function. The ethical concern is unjustifiable disadvantage; every algorithm embeds values, so the goal is avoiding problematic bias, not all bias.
\item \textbf{Why B is wrong:} The fee example shows unbiased is not the same as ethical
\item \textbf{Why C is wrong:} The module says no algorithm is objective; each reflects embedded values
\item \textbf{Why D is wrong:} Problematic biases often occur unintentionally due to complexity
\end{itemize}

\textcolor{myred}{\textbf{EXAM TRAP:}} Not all bias is problematic, and not all unbiased AI is ethical!
\end{frame}

\begin{frame}[shrink=6]{Q21: Problem Specification Bias - Advanced}
\fontsize{6.5}{7.5}\selectfont
\textbf{QUESTION: A bank wants zero lending risk and its algorithm ends up serving only affluent applicants. Which diagnosis and remedy fit best?}

\vspace{0.1cm}
\textbf{OPTIONS:}
\begin{enumerate}[A.]
\item Historical bias; add more decades of lending data
\item Bias in problem specification: the goal is operationalized in a way that misses the real-world objective, so the target should be modified to accommodate reasonable risk
\item Sampling bias; oversample low-income applicants only
\item Deployment bias; retrain loan officers
\end{enumerate}

\vspace{0.1cm}
\textcolor{myblue}{\textbf{ANSWER: B}}

\vspace{0.1cm}
\textbf{DETAILED EXPLANATION:}
\begin{itemize}
\item \textbf{Why B is correct:} The flaw lies in the goal itself, not the data or its use. Zero-risk targets exclude everyone who cannot guarantee repayment.
\item \textbf{Why A is wrong:} More historical data would reinforce the same target
\item \textbf{Why C is wrong:} The issue is the objective, not the sample
\item \textbf{Why D is wrong:} The problem exists before deployment
\end{itemize}

\textcolor{myred}{\textbf{EXAM TRAP:}} Bias can be born in the GOAL: fix the target, not just the data!
\end{frame}

\begin{frame}[shrink=6]{Q22: Proxy Targets: The Healthcare Example - Advanced}
\fontsize{6.5}{7.5}\selectfont
\textbf{QUESTION: An algorithm meant to identify the sickest patients used healthcare expenditure as a proxy and favored wealthier patients. Which mechanism explains this?}

\vspace{0.1cm}
\textbf{OPTIONS:}
\begin{enumerate}[A.]
\item Selective labels: outcomes were unobserved for patients not treated
\item Automation bias: clinicians over-trusted the tool
\item Sampling bias: too few patients were sampled
\item The proxy diverges from the real objective (illness) because spending reflects access to care and wealth, so optimizing it selects the richest rather than the sickest
\end{enumerate}

\vspace{0.1cm}
\textcolor{myblue}{\textbf{ANSWER: D}}

\vspace{0.1cm}
\textbf{DETAILED EXPLANATION:}
\begin{itemize}
\item \textbf{Why D is correct:} This is problem-specification bias through a flawed proxy: the measurable stand-in encodes unequal spending patterns.
\item \textbf{Why A is wrong:} Selective labels concerns missing counterfactual outcomes, not the choice of target
\item \textbf{Why B is wrong:} Automation bias is a deployment issue about deference
\item \textbf{Why C is wrong:} The failure exists even if the sample is large and random
\end{itemize}

\textcolor{myred}{\textbf{EXAM TRAP:}} Cost is not need: a bad PROXY target creates biased outcomes!
\end{frame}

\begin{frame}[shrink=6]{Q23: Historical Bias Beyond Discrimination - Advanced}
\fontsize{6.5}{7.5}\selectfont
\textbf{QUESTION: A model trained on 2010-2019 data forecasts office vacancy and commuter rail volume for 2025-2030. Which statement fits the module's discussion of historical data?}

\vspace{0.1cm}
\textbf{OPTIONS:}
\begin{enumerate}[A.]
\item Historical bias exists only when data reflect sex or race disparities
\item Historical data are reliable if they are large and accurate for their period
\item Historical data can mislead even without demographic bias: patterns that held in the past may not hold, producing inaccurate or malfunctioning models
\item The problem is deployment bias created by human deference
\end{enumerate}

\vspace{0.1cm}
\textcolor{myblue}{\textbf{ANSWER: C}}

\vspace{0.1cm}
\textbf{DETAILED EXPLANATION:}
\begin{itemize}
\item \textbf{Why C is correct:} The module extends historical bias beyond sex and race to mismatch between past and future real-world trends, illustrated by pre-pandemic versus later patterns.
\item \textbf{Why A is wrong:} The module explicitly says it also leads to inaccurate models more broadly
\item \textbf{Why B is wrong:} Accuracy for the past period does not guarantee validity for a changed future
\item \textbf{Why D is wrong:} No human deference is involved in this example
\end{itemize}

\textcolor{myred}{\textbf{EXAM TRAP:}} Historical bias includes DATA THAT NO LONGER MATCHES THE WORLD!
\end{frame}

\begin{frame}[shrink=6]{Q24: Selective Labels vs Sampling Bias - Advanced}
\fontsize{6.5}{7.5}\selectfont
\textbf{QUESTION: Which statement correctly distinguishes selective labels from sampling bias?}

\vspace{0.1cm}
\textbf{OPTIONS:}
\begin{enumerate}[A.]
\item Selective labels: outcomes are observed only for people selected by past decisions, so counterfactuals are missing; sampling bias: the sample is not random, so trends may not generalize
\item They are the same problem described in two fields
\item Selective labels: data drawn from too few people; sampling bias: missing outcomes for those denied
\item Selective labels arise in deployment; sampling bias arises in problem specification
\end{enumerate}

\vspace{0.1cm}
\textcolor{myblue}{\textbf{ANSWER: A}}

\vspace{0.1cm}
\textbf{DETAILED EXPLANATION:}
\begin{itemize}
\item \textbf{Why A is correct:} A lender sees repayment only for approved borrowers, so it never learns whether rejected ones would have repaid. Sampling bias, such as drug data mostly from young men, limits generalization.
\item \textbf{Why B is wrong:} The module calls them related but distinct
\item \textbf{Why C is wrong:} This swaps the definitions
\item \textbf{Why D is wrong:} Both are data biases, not deployment or problem-specification biases
\end{itemize}

\textcolor{myred}{\textbf{EXAM TRAP:}} Selective labels = missing COUNTERFACTUAL outcomes; sampling bias = NON-RANDOM sample!
\end{frame}

\begin{frame}[shrink=6]{Q25: Popularity Feedback Loops - Advanced}
\fontsize{6.5}{7.5}\selectfont
\textbf{QUESTION: A search engine for financial products favors popular items, giving them more exposure and thereby more popularity, regardless of quality. Which category is this?}

\vspace{0.1cm}
\textbf{OPTIONS:}
\begin{enumerate}[A.]
\item Bias in problem specification
\item Bias in modeling, validation, and algorithm design, which can arise even when specification and data are sound
\item Selective labels bias
\item Bias in deployment caused by human deference
\end{enumerate}

\vspace{0.1cm}
\textcolor{myblue}{\textbf{ANSWER: B}}

\vspace{0.1cm}
\textbf{DETAILED EXPLANATION:}
\begin{itemize}
\item \textbf{Why B is correct:} Choices about optimization functions, subgroups, and presentation can introduce bias. Here the design creates a self-reinforcing exposure loop.
\item \textbf{Why A is wrong:} The goal (helping users choose) is not the flawed part; the design is
\item \textbf{Why C is wrong:} Selective labels concerns missing outcomes for those not selected
\item \textbf{Why D is wrong:} No human deference is described
\end{itemize}

\textcolor{myred}{\textbf{EXAM TRAP:}} Design choices can create FEEDBACK LOOPS even with good goals and data!
\end{frame}

\begin{frame}[shrink=6]{Q26: Deployment Bias and Deference - Advanced}
\fontsize{6.5}{7.5}\selectfont
\textbf{QUESTION: A bank warns employees not to rely solely on a loan-risk algorithm still under testing, yet many follow it uncritically. Which statement is accurate?}

\vspace{0.1cm}
\textbf{OPTIONS:}
\begin{enumerate}[A.]
\item Not a bias issue, because the algorithm itself is unbiased
\item It is historical bias, fixable by retraining on more data
\item It is caused by algorithms being more accurate than people
\item This is deployment bias: an unbiased algorithm can still yield problematic outcomes because deference shields people from blame, so implementation changes behavior
\end{enumerate}

\vspace{0.1cm}
\textcolor{myblue}{\textbf{ANSWER: D}}

\vspace{0.1cm}
\textbf{DETAILED EXPLANATION:}
\begin{itemize}
\item \textbf{Why D is correct:} The module lists hypotheses: convenience, perceived objectivity, self-doubt, and sharing of responsibility. Deferring lowers the risk of harsher judgment when errors occur.
\item \textbf{Why A is wrong:} The module says biases can emerge at deployment despite an unbiased algorithm
\item \textbf{Why B is wrong:} Retraining does not change the incentive to defer
\item \textbf{Why C is wrong:} The account rests on incentives, not on proven accuracy
\end{itemize}

\textcolor{myred}{\textbf{EXAM TRAP:}} Deployment bias: unintended INCENTIVE EFFECTS make people defer to the algorithm!
\end{frame}

\begin{frame}[shrink=6]{Q27: The Fairness Impossibility Result - Advanced}
\fontsize{6.5}{7.5}\selectfont
\textbf{QUESTION: According to the module, which set of fairness criteria cannot be satisfied simultaneously when base rates differ across groups?}

\vspace{0.1cm}
\textbf{OPTIONS:}
\begin{enumerate}[A.]
\item Procedural fairness, distributive justice, and restorative justice
\item Group fairness and privacy, unless the data are synthetic
\item Individual fairness and explainability
\item Demographic parity, predictive rate parity, and equal opportunity
\end{enumerate}

\vspace{0.1cm}
\textcolor{myblue}{\textbf{ANSWER: D}}

\vspace{0.1cm}
\textbf{DETAILED EXPLANATION:}
\begin{itemize}
\item \textbf{Why D is correct:} Automating fairness becomes feasible only when base rates are equal, which is rare. Cited work includes Kleinberg, Mullainathan, and Raghavan (2016) and Chouldechova (2017).
\item \textbf{Why A is wrong:} Those are conceptions of justice, not statistical criteria
\item \textbf{Why B is wrong:} Privacy and synthetic data are unrelated to the result
\item \textbf{Why C is wrong:} The result concerns three statistical fairness criteria
\end{itemize}

\textcolor{myred}{\textbf{EXAM TRAP:}} Unequal base rates = cannot satisfy demographic parity, predictive parity, and equal opportunity at once!
\end{frame}

\begin{frame}[shrink=6]{Q28: Can Fairness Auditing Be Automated? - Advanced}
\fontsize{6.5}{7.5}\selectfont
\textbf{QUESTION: The module argues that AI cannot solve the fairness problems it creates. Which chain of reasoning is correct?}

\vspace{0.1cm}
\textbf{OPTIONS:}
\begin{enumerate}[A.]
\item Fairness cannot be automated because it involves moral trade-offs, so auditing for fairness cannot be fully automated either; tools help but rarely solve bias on their own
\item Fairness is mathematical, but tools are too slow
\item Auditing tools are unreliable, so audits should be abandoned
\item AI-based auditors avoid the problem because they lack values
\end{enumerate}

\vspace{0.1cm}
\textcolor{myblue}{\textbf{ANSWER: A}}

\vspace{0.1cm}
\textbf{DETAILED EXPLANATION:}
\begin{itemize}
\item \textbf{Why A is correct:} Fairness is a moral or ethical judgment, not a mathematical one, and involves imperfect compromises that change with circumstances. Audits also need a fresh, diverse human look.
\item \textbf{Why B is wrong:} The module says fairness is not purely mathematical
\item \textbf{Why C is wrong:} Auditing is called likely the best way to find and correct bias
\item \textbf{Why D is wrong:} The module says every algorithm reflects values
\end{itemize}

\textcolor{myred}{\textbf{EXAM TRAP:}} Fairness = moral judgment, so tools help but the audit cannot be fully automated!
\end{frame}

\begin{frame}[shrink=6]{Q29: Risk Scores and Unequal Base Rates - Advanced}
\fontsize{6.5}{7.5}\selectfont
\textbf{QUESTION: Most people committing a certain crime are men. An AI scores a specific man and woman. Which conclusion follows from the module?}

\vspace{0.1cm}
\textbf{OPTIONS:}
\begin{enumerate}[A.]
\item Equal scores are always fair because individuals should be treated alike
\item A higher score for the man is always fair because it is statistically justified
\item Whatever scoring rule is chosen, one person can claim unfair treatment, showing why fairness is a moral trade-off rather than a computation
\item Fairness disputes vanish once the model is accurate
\end{enumerate}

\vspace{0.1cm}
\textcolor{myblue}{\textbf{ANSWER: C}}

\vspace{0.1cm}
\textbf{DETAILED EXPLANATION:}
\begin{itemize}
\item \textbf{Why C is correct:} A higher score for the man invites his complaint; a similar score invites the woman's complaint because her group's base rate is lower.
\item \textbf{Why A is wrong:} Equal scores draw the woman's objection
\item \textbf{Why B is wrong:} A higher score draws the man's objection
\item \textbf{Why D is wrong:} Accuracy does not resolve conflicts among fairness criteria
\end{itemize}

\textcolor{myred}{\textbf{EXAM TRAP:}} With unequal base rates, every choice leaves someone with a fairness complaint!
\end{frame}

\begin{frame}[shrink=6]{Q30: Procedural Fairness and Redress - Advanced}
\fontsize{6.5}{7.5}\selectfont
\textbf{QUESTION: In the justice-system analogy, why can errors be justifiable under procedural fairness? Which AI analogue follows?}

\vspace{0.1cm}
\textbf{OPTIONS:}
\begin{enumerate}[A.]
\item Outcomes are always correct when the process is fair
\item Only outcome fairness matters, so process is secondary
\item Errors are never justifiable, so committees are unnecessary
\item A fair, impartial process with means of redress makes errors defensible and correctable; firms need structures such as ethics committees weighing all three frameworks
\end{enumerate}

\vspace{0.1cm}
\textcolor{myblue}{\textbf{ANSWER: D}}

\vspace{0.1cm}
\textbf{DETAILED EXPLANATION:}
\begin{itemize}
\item \textbf{Why D is correct:} Mistakes can occur even with fair procedures, but evidence-based process and mechanisms like new trials allow redress. The challenge is to build corporate structures that deliver both procedural and outcome fairness.
\item \textbf{Why A is wrong:} Mistakes still happen with fair procedures
\item \textbf{Why B is wrong:} The module says fairness is not only about outcome
\item \textbf{Why C is wrong:} The module holds that committees can support both kinds of fairness
\end{itemize}

\textcolor{myred}{\textbf{EXAM TRAP:}} Procedural fairness = impartial process plus ability to RIGHT WRONGS!
\end{frame}

\begin{frame}[shrink=6]{Q31: Synthetic Data Trade-offs - Advanced}
\fontsize{6.5}{7.5}\selectfont
\textbf{QUESTION: Which statement about synthetic data as a bias-mitigation tool is accurate?}

\vspace{0.1cm}
\textbf{OPTIONS:}
\begin{enumerate}[A.]
\item It always outperforms real data, so validation is unnecessary
\item It avoids privacy risk to real data subjects and can be designed to remove bias, but it can be less precise (especially visual data) and differ from real data in non-obvious ways
\item It is bias-free automatically because it is fabricated
\item It eliminates privacy and quality concerns entirely
\end{enumerate}

\vspace{0.1cm}
\textcolor{myblue}{\textbf{ANSWER: B}}

\vspace{0.1cm}
\textbf{DETAILED EXPLANATION:}
\begin{itemize}
\item \textbf{Why B is correct:} Synthetic data are fabricated datasets. The advantage comes from design choices, and the disadvantage is fidelity.
\item \textbf{Why A is wrong:} The module says precision can be lower
\item \textbf{Why C is wrong:} Bias must be removed by design, not automatically
\item \textbf{Why D is wrong:} Quality concerns remain
\end{itemize}

\textcolor{myred}{\textbf{EXAM TRAP:}} Synthetic data: privacy and bias-control benefits, but FIDELITY risk!
\end{frame}

\begin{frame}[shrink=6]{Q32: Self-Driving Trade-offs - Advanced}
\fontsize{6.5}{7.5}\selectfont
\textbf{QUESTION: A self-driving system is tuned to be especially sensitive to children, slightly reducing adult detection accuracy. Which pairing is correct?}

\vspace{0.1cm}
\textbf{OPTIONS:}
\begin{enumerate}[A.]
\item The aim is avoiding problematic biases, not all bias; a deontological check is that protected categories are not disfavored (for example, better detection of larger men) and safety minimums hold
\item Any accuracy trade-off is unethical because bias must be eliminated
\item A consequentialist check: ensure no protected category is disfavored
\item A virtue check: compute accident risk for each version
\end{enumerate}

\vspace{0.1cm}
\textcolor{myblue}{\textbf{ANSWER: A}}

\vspace{0.1cm}
\textbf{DETAILED EXPLANATION:}
\begin{itemize}
\item \textbf{Why A is correct:} Trade-offs must be justifiable to stakeholders. Consequentialists compute accident risk, deontologists guard protected categories and minimums, and virtue ethicists focus on responsible processes.
\item \textbf{Why B is wrong:} The module says the goal is not to avoid bias in general
\item \textbf{Why C is wrong:} Protected categories and safety minimums are deontological
\item \textbf{Why D is wrong:} Risk calculation is consequentialist
\end{itemize}

\textcolor{myred}{\textbf{EXAM TRAP:}} Avoid PROBLEMATIC bias; consequences = risk math, deontology = protected groups and minimums, virtue = responsible process!
\end{frame}

\begin{frame}[shrink=6]{Q33: Bias vs Legal Discrimination - Advanced}
\fontsize{6.5}{7.5}\selectfont
\textbf{QUESTION: An algorithm disfavors applicants on a characteristic that is irrelevant to the task but not legally protected in the jurisdiction. What is the best assessment?}

\vspace{0.1cm}
\textbf{OPTIONS:}
\begin{enumerate}[A.]
\item It is neither, since only protected classes matter
\item It is automatically illegal discrimination everywhere
\item It is a problematic bias ethically, because it deviates from function based on a morally irrelevant trait, although whether it is legal discrimination depends on jurisdiction and protected classes
\item It is acceptable because it is statistically accurate
\end{enumerate}

\vspace{0.1cm}
\textcolor{myblue}{\textbf{ANSWER: C}}

\vspace{0.1cm}
\textbf{DETAILED EXPLANATION:}
\begin{itemize}
\item \textbf{Why C is correct:} Legal discrimination is likely when disfavoring people by race, sex, ethnicity, age, or another protected class. Ethical fairness concerns irrelevant characteristics more broadly.
\item \textbf{Why A is wrong:} Ethics reaches beyond legal minimums
\item \textbf{Why B is wrong:} Discrimination law varies by jurisdiction
\item \textbf{Why D is wrong:} Statistical unbiasedness does not guarantee ethical acceptability
\end{itemize}

\textcolor{myred}{\textbf{EXAM TRAP:}} Ethics goes beyond law: legal discrimination varies by JURISDICTION!
\end{frame}

% ===== PRIVACY =====

\begin{frame}[shrink=6]{Q34: Access-Based Privacy - Advanced}
\fontsize{6.5}{7.5}\selectfont
\textbf{QUESTION: The module defines privacy as others having no access to a personal data point. Under this definition, which event is a privacy loss?}

\vspace{0.1cm}
\textbf{OPTIONS:}
\begin{enumerate}[A.]
\item Only a public leak of the data
\item Only use of the data for discrimination
\item Only access that violates a statute
\item An employer gains access to a personal data point, even if it never misuses it
\end{enumerate}

\vspace{0.1cm}
\textcolor{myblue}{\textbf{ANSWER: D}}

\vspace{0.1cm}
\textbf{DETAILED EXPLANATION:}
\begin{itemize}
\item \textbf{Why D is correct:} Privacy is proportional to lack of access, and it protects from abuses of power because more knowledge makes interference easier.
\item \textbf{Why A is wrong:} Loss occurs on access, not only on publication
\item \textbf{Why B is wrong:} Misuse is a possible harm, not the definition
\item \textbf{Why C is wrong:} The moral right holds whether or not the law recognizes it
\end{itemize}

\textcolor{myred}{\textbf{EXAM TRAP:}} Privacy = lack of ACCESS; loss does not require misuse!
\end{frame}

\begin{frame}[shrink=6]{Q35: Personal Data as a Toxic Asset - Advanced}
\fontsize{6.5}{7.5}\selectfont
\textbf{QUESTION: Which argument best supports data minimization on the module's terms?}

\vspace{0.1cm}
\textbf{OPTIONS:}
\begin{enumerate}[A.]
\item Minimization mainly lowers storage costs
\item Minimization applies only to regulated data
\item Data are cheap to mine and valuable but sensitive and hard to keep safe; each data point is a potential lawsuit or fine, and risk grows with amount, storage time, analysis, and sharing
\item Data are risk-free if encrypted
\end{enumerate}

\vspace{0.1cm}
\textcolor{myblue}{\textbf{ANSWER: C}}

\vspace{0.1cm}
\textbf{DETAILED EXPLANATION:}
\begin{itemize}
\item \textbf{Why C is correct:} Data are a toxic yet valuable asset. Collect only when the benefit to data subjects outweighs disadvantages.
\item \textbf{Why A is wrong:} Cost is secondary to risk and rights
\item \textbf{Why B is wrong:} The principle is ethical and general
\item \textbf{Why D is wrong:} Security cannot eliminate risk given attacker advantages
\end{itemize}

\textcolor{myred}{\textbf{EXAM TRAP:}} More data collected, kept, analyzed, and shared = more risk!
\end{frame}

\begin{frame}[shrink=6]{Q36: Attacker-Defender Asymmetry - Advanced}
\fontsize{6.5}{7.5}\selectfont
\textbf{QUESTION: In cyberspace, defenders are at a disadvantage. What follows for a firm holding excess personal data?}

\vspace{0.1cm}
\textbf{OPTIONS:}
\begin{enumerate}[A.]
\item Attacks are avoidable with enough encryption
\item Only breaches that occur are ethically relevant
\item If an organization cannot keep data safe it puts itself at risk by collecting it, and imposing unjustified breach risk on data subjects can itself be a wrong
\item Risk is borne only by the firm
\end{enumerate}

\vspace{0.1cm}
\textcolor{myblue}{\textbf{ANSWER: C}}

\vspace{0.1cm}
\textbf{DETAILED EXPLANATION:}
\begin{itemize}
\item \textbf{Why C is correct:} Attackers choose time and method; defenders must guard against all. Risk imposition can amount to a wrong even if the attempt fails.
\item \textbf{Why A is wrong:} A resourced, motivated attacker will eventually succeed
\item \textbf{Why B is wrong:} Imposing risk can be wrong without a completed harm
\item \textbf{Why D is wrong:} Risk also falls on data subjects and society
\end{itemize}

\textcolor{myred}{\textbf{EXAM TRAP:}} Defenders must stop EVERY attack; collecting unnecessary data imposes risk!
\end{frame}

\begin{frame}[shrink=6]{Q37: Contextual Integrity - Advanced}
\fontsize{6.5}{7.5}\selectfont
\textbf{QUESTION: A bank stores customer data with strong encryption but sells it to a marketing company. Which principle is violated most directly?}

\vspace{0.1cm}
\textbf{OPTIONS:}
\begin{enumerate}[A.]
\item Data deletion: retention beyond need
\item Right to be forgotten: removal from search results
\item Data security: the encryption was inadequate
\item Contextual integrity: data given in one context should not be moved to another that violates its norms
\end{enumerate}

\vspace{0.1cm}
\textcolor{myblue}{\textbf{ANSWER: D}}

\vspace{0.1cm}
\textbf{DETAILED EXPLANATION:}
\begin{itemize}
\item \textbf{Why D is correct:} The customer provided data for a transaction. Moving it to marketing violates the norms of that context even when the system is secure.
\item \textbf{Why A is wrong:} Deletion concerns timing, not repurposing
\item \textbf{Why B is wrong:} That right concerns inadequate, irrelevant, or excessive information in searches and directories
\item \textbf{Why C is wrong:} The scenario stipulates strong encryption
\end{itemize}

\textcolor{myred}{\textbf{EXAM TRAP:}} Secure is not the same as appropriate: contextual integrity governs USE across contexts!
\end{frame}

\begin{frame}[shrink=6]{Q38: Right to Be Forgotten - Advanced}
\fontsize{6.5}{7.5}\selectfont
\textbf{QUESTION: When does the right to be forgotten apply, according to the module?}

\vspace{0.1cm}
\textbf{OPTIONS:}
\begin{enumerate}[A.]
\item Whenever a person finds information embarrassing
\item When information is inadequate (for example incorrect), irrelevant or no longer relevant, or excessive relative to its purposes, making it less accessible in searches and directories
\item Only for criminal records
\item Only when data was obtained unlawfully
\end{enumerate}

\vspace{0.1cm}
\textcolor{myblue}{\textbf{ANSWER: B}}

\vspace{0.1cm}
\textbf{DETAILED EXPLANATION:}
\begin{itemize}
\item \textbf{Why B is correct:} It concerns removal from internet searches and other directories under specified conditions.
\item \textbf{Why A is wrong:} Embarrassment alone is not the stated criterion
\item \textbf{Why C is wrong:} It is not limited to criminal records
\item \textbf{Why D is wrong:} Lawful data may still become irrelevant
\end{itemize}

\textcolor{myred}{\textbf{EXAM TRAP:}} Right to be forgotten = inadequate, irrelevant, or excessive data!
\end{frame}

\begin{frame}[shrink=6]{Q39: Control Over Data - Advanced}
\fontsize{6.5}{7.5}\selectfont
\textbf{QUESTION: Which element of control over data goes beyond the raw personal data supplied by individuals?}

\vspace{0.1cm}
\textbf{OPTIONS:}
\begin{enumerate}[A.]
\item Ability to withdraw consent
\item Ability to delete their data
\item Access to the inferences made about them
\item Ability to transfer their data
\end{enumerate}

\vspace{0.1cm}
\textcolor{myblue}{\textbf{ANSWER: C}}

\vspace{0.1cm}
\textbf{DETAILED EXPLANATION:}
\begin{itemize}
\item \textbf{Why C is correct:} Control includes denying or withdrawing consent, accessing data and inferences, transfer, and deletion. Inferences are derived rather than supplied.
\item \textbf{Why A is wrong:} Withdrawal concerns consent to raw collection and processing
\item \textbf{Why B is wrong:} Deletion concerns stored data
\item \textbf{Why D is wrong:} Portability concerns moving data
\end{itemize}

\textcolor{myred}{\textbf{EXAM TRAP:}} Control includes access to INFERENCES, not only raw data!
\end{frame}

% ===== GOVERNANCE =====

\begin{frame}[shrink=6]{Q40: Backpropagation and Interpretability - Advanced}
\fontsize{6.5}{7.5}\selectfont
\textbf{QUESTION: Why does training by backpropagation contribute to the interpretability problem?}

\vspace{0.1cm}
\textbf{OPTIONS:}
\begin{enumerate}[A.]
\item It randomizes weights so outputs are unpredictable
\item It adjusts weights to minimize output error, improving performance without revealing how decisions are reached, and scale (millions or billions of parameters) makes tracing difficult
\item It keeps model weights secret from developers
\item It prevents networks from using proxies
\end{enumerate}

\vspace{0.1cm}
\textcolor{myblue}{\textbf{ANSWER: B}}

\vspace{0.1cm}
\textbf{DETAILED EXPLANATION:}
\begin{itemize}
\item \textbf{Why B is correct:} Effectiveness at reducing error does not provide insight into reasoning. Auditing methods are being developed to look into the black box.
\item \textbf{Why A is wrong:} Backpropagation is systematic error minimization
\item \textbf{Why C is wrong:} The problem persists even when weights are visible
\item \textbf{Why D is wrong:} Proxies are common, and unknown proxies hamper governance
\end{itemize}

\textcolor{myred}{\textbf{EXAM TRAP:}} Backpropagation optimizes error, NOT understanding!
\end{frame}

\begin{frame}[shrink=6]{Q41: Unpredictability, RCTs, and Audits - Advanced}
\fontsize{6.5}{7.5}\selectfont
\textbf{QUESTION: Why does the module suggest both randomized controlled trials and periodic audits for unpredictable AI?}

\vspace{0.1cm}
\textbf{OPTIONS:}
\begin{enumerate}[A.]
\item Emergent behavior and unforeseen human interactions mean pre-deployment trials alone cannot capture all failure modes, so audits are a complementary check on safety and accuracy
\item RCTs eliminate unpredictability, so audits are redundant
\item Audits replace RCTs because trials are unethical
\item Unpredictability arises only from user misuse
\end{enumerate}

\vspace{0.1cm}
\textcolor{myblue}{\textbf{ANSWER: A}}

\vspace{0.1cm}
\textbf{DETAILED EXPLANATION:}
\begin{itemize}
\item \textbf{Why A is correct:} Neural networks exhibit emergent properties not coded explicitly and can surprise experts. Misuse is also hard to predict.
\item \textbf{Why B is wrong:} RCTs minimize but do not eliminate risk
\item \textbf{Why C is wrong:} The module calls audits complementary to RCTs
\item \textbf{Why D is wrong:} Unpredictability also arises from emergent capabilities and errors
\end{itemize}

\textcolor{myred}{\textbf{EXAM TRAP:}} Emergent behavior: RCTs AND periodic audits!
\end{frame}

\begin{frame}[shrink=6]{Q42: LLMs and Truth-Tracking - Advanced}
\fontsize{6.5}{7.5}\selectfont
\textbf{QUESTION: Why are LLM outputs risky even when fluent?}

\vspace{0.1cm}
\textbf{OPTIONS:}
\begin{enumerate}[A.]
\item They deliberately deceive users
\item They rely on verified knowledge bases
\item They are unregulated only in the EU
\item They make statistical inferences designed to be plausible, and plausible is not the same as truthful, creating risks of misinformation, physical safety issues, and libel
\end{enumerate}

\vspace{0.1cm}
\textcolor{myblue}{\textbf{ANSWER: D}}

\vspace{0.1cm}
\textbf{DETAILED EXPLANATION:}
\begin{itemize}
\item \textbf{Why D is correct:} LLMs are not based on an understanding of truth or knowledge of the world, and false responses can still be convincing.
\item \textbf{Why A is wrong:} The issue is design, not intent
\item \textbf{Why B is wrong:} The module says they are not truth-tracking
\item \textbf{Why C is wrong:} The concern is global
\end{itemize}

\textcolor{myred}{\textbf{EXAM TRAP:}} Plausible does not mean true!
\end{frame}

\begin{frame}[shrink=6]{Q43: LLMs and GDPR Rights - Advanced}
\fontsize{6.5}{7.5}\selectfont
\textbf{QUESTION: Why does the module question whether LLM providers can comply with the GDPR?}

\vspace{0.1cm}
\textbf{OPTIONS:}
\begin{enumerate}[A.]
\item The GDPR does not apply to scraped data
\item Data subjects may ask what data a company holds, modify it, and delete it, and it is unclear providers can meet such requests for data scraped and absorbed into models
\item LLMs never process personal data
\item The GDPR bans language models
\end{enumerate}

\vspace{0.1cm}
\textcolor{myblue}{\textbf{ANSWER: B}}

\vspace{0.1cm}
\textbf{DETAILED EXPLANATION:}
\begin{itemize}
\item \textbf{Why B is correct:} Scraping raises privacy and copyright concerns, and lawsuits were pending at the time of writing.
\item \textbf{Why A is wrong:} The GDPR applies to personal data generally
\item \textbf{Why C is wrong:} The module raises scraped personal data as an issue
\item \textbf{Why D is wrong:} The GDPR is not a ban on any technology
\end{itemize}

\textcolor{myred}{\textbf{EXAM TRAP:}} GDPR rights (access, correct, delete) are hard to honor for scraped LLM training data!
\end{frame}

% ===== REGULATION =====

\begin{frame}[shrink=6]{Q44: GDPR Reach and Breach Notification - Advanced}
\fontsize{6.5}{7.5}\selectfont
\textbf{QUESTION: A US company with no EU office offers a service to people in the EEA and suffers a breach. What applies?}

\vspace{0.1cm}
\textbf{OPTIONS:}
\begin{enumerate}[A.]
\item The GDPR does not apply outside the EEA
\item Notification is due within 45 days
\item Notification is due only if individuals sue
\item The GDPR applies because the company offers services to EEA data subjects; the controller must notify the supervisory authority within 72 hours
\end{enumerate}

\vspace{0.1cm}
\textcolor{myblue}{\textbf{ANSWER: D}}

\vspace{0.1cm}
\textbf{DETAILED EXPLANATION:}
\begin{itemize}
\item \textbf{Why D is correct:} Extraterritorial reach has made the GDPR globally influential. Controllers must notify within 72 hours.
\item \textbf{Why A is wrong:} It applies based on offering services to EEA data subjects
\item \textbf{Why B is wrong:} 45 days is the appendix figure for the Illinois BIPA
\item \textbf{Why C is wrong:} The duty is legal and does not depend on suits
\end{itemize}

\textcolor{myred}{\textbf{EXAM TRAP:}} GDPR: extraterritorial, 72-hour breach notice!
\end{frame}

\begin{frame}[shrink=6]{Q45: Legitimate Interest - Advanced}
\fontsize{6.5}{7.5}\selectfont
\textbf{QUESTION: Under the GDPR, when may the right to object to processing not apply?}

\vspace{0.1cm}
\textbf{OPTIONS:}
\begin{enumerate}[A.]
\item Whenever the controller is outside the EEA
\item Whenever data are anonymized
\item When there is a legitimate interest to carry out the service, whose scope is debated and being clarified by rulings
\item Never; the right is absolute
\end{enumerate}

\vspace{0.1cm}
\textcolor{myblue}{\textbf{ANSWER: C}}

\vspace{0.1cm}
\textbf{DETAILED EXPLANATION:}
\begin{itemize}
\item \textbf{Why C is correct:} The right to object covers marketing or purposes unrelated to the service, but does not apply where legitimate interest exists.
\item \textbf{Why A is wrong:} Location does not remove GDPR duties when serving EEA data subjects
\item \textbf{Why B is wrong:} Anonymization is unrelated to this exception
\item \textbf{Why D is wrong:} The module lists legitimate interest as a limit
\end{itemize}

\textcolor{myred}{\textbf{EXAM TRAP:}} Right to object has a LEGITIMATE-INTEREST limit!
\end{frame}

\begin{frame}[shrink=6]{Q46: EU AI Act: Prohibited vs Regulated - Advanced}
\fontsize{6.5}{7.5}\selectfont
\textbf{QUESTION: Which system is among those prohibited under the EU AI Act, as opposed to merely regulated?}

\vspace{0.1cm}
\textbf{OPTIONS:}
\begin{enumerate}[A.]
\item An AI system that screens job applicants
\item A chatbot that interacts with customers
\item Untargeted scraping of facial images from the internet or CCTV to build facial recognition databases
\item A generator of deepfake content
\end{enumerate}

\vspace{0.1cm}
\textcolor{myblue}{\textbf{ANSWER: C}}

\vspace{0.1cm}
\textbf{DETAILED EXPLANATION:}
\begin{itemize}
\item \textbf{Why C is correct:} Prohibited systems include biometric categorization using sensitive traits, untargeted facial scraping, workplace and educational emotion recognition, social scoring, manipulative systems, and exploiting vulnerabilities.
\item \textbf{Why A is wrong:} Job screening carries obligations such as preventing undue harm to opportunities
\item \textbf{Why B is wrong:} Chatbots require notification
\item \textbf{Why D is wrong:} Deepfakes require labeling
\end{itemize}

\textcolor{myred}{\textbf{EXAM TRAP:}} Facial scraping banned; job-screening AI, chatbots and deepfakes regulated!
\end{frame}

\begin{frame}[shrink=6]{Q47: EU AI Act Penalties and Timeline - Advanced}
\fontsize{6.5}{7.5}\selectfont
\textbf{QUESTION: Which pairing matches the module's account of AI Act penalties and timing?}

\vspace{0.1cm}
\textbf{OPTIONS:}
\begin{enumerate}[A.]
\item Fines range from EUR 35 million or 7\% of global turnover down to EUR 7.5 million or 1.5\%; prohibitions applied from February 2, 2025 and most provisions from August 2, 2026
\item Fines are capped at EUR 1 million; effective January 1, 2020
\item Fines equal 72 hours of turnover; all provisions applied on August 1, 2024
\item There are no fines; the AI Office only issues guidance
\end{enumerate}

\vspace{0.1cm}
\textcolor{myblue}{\textbf{ANSWER: A}}

\vspace{0.1cm}
\textbf{DETAILED EXPLANATION:}
\begin{itemize}
\item \textbf{Why A is correct:} The Act entered into force August 1, 2024 and created the European AI Office for compliance and enforcement.
\item \textbf{Why B is wrong:} This misstates both figures and dates
\item \textbf{Why C is wrong:} 72 hours is the GDPR breach window
\item \textbf{Why D is wrong:} The AI Office enforces binding rules
\end{itemize}

\textcolor{myred}{\textbf{EXAM TRAP:}} AI Act: up to EUR 35M or 7\%; bans Feb 2, 2025; most provisions Aug 2, 2026!
\end{frame}

\begin{frame}[shrink=6]{Q48: High-Risk Obligations - Advanced}
\fontsize{6.5}{7.5}\selectfont
\textbf{QUESTION: Which allocation of high-risk AI obligations is accurate?}

\vspace{0.1cm}
\textbf{OPTIONS:}
\begin{enumerate}[A.]
\item Only chatbot operators must assess rights impact
\item Providers need not keep records if they publish source code
\item Only public agencies must document training
\item Banks, insurers, and deployers of high-risk AI do rights impact assessments; providers keep records on datasets, training methods, and oversight measures
\end{enumerate}

\vspace{0.1cm}
\textcolor{myblue}{\textbf{ANSWER: D}}

\vspace{0.1cm}
\textbf{DETAILED EXPLANATION:}
\begin{itemize}
\item \textbf{Why D is correct:} The Act imposes assessments on essential-service organizations and record-keeping duties on providers.
\item \textbf{Why A is wrong:} Chatbots have notification duties instead
\item \textbf{Why B is wrong:} No source-code exemption is described
\item \textbf{Why C is wrong:} The duties are not limited to public agencies
\end{itemize}

\textcolor{myred}{\textbf{EXAM TRAP:}} High-risk: impact assessments plus thorough records!
\end{frame}

\begin{frame}[shrink=6]{Q49: DMA vs DSA - Advanced}
\fontsize{6.5}{7.5}\selectfont
\textbf{QUESTION: A platform ranks its own products above rivals; another very large platform must let users opt out of personalized recommendations. Which laws apply?}

\vspace{0.1cm}
\textbf{OPTIONS:}
\begin{enumerate}[A.]
\item Both are DMA obligations
\item The DMA bans self-preferencing by gatekeepers; the DSA requires very large platforms to allow opt-out of personalized content
\item The GDPR bans self-preferencing; the AI Act covers opt-outs
\item Both are DSA obligations on content moderation only
\end{enumerate}

\vspace{0.1cm}
\textcolor{myblue}{\textbf{ANSWER: B}}

\vspace{0.1cm}
\textbf{DETAILED EXPLANATION:}
\begin{itemize}
\item \textbf{Why B is correct:} The DMA targets gatekeepers (self-preferencing, data access, no tying, messaging interoperability). The DSA covers illegal content, transparency, disinformation, and algorithmic personalization.
\item \textbf{Why A is wrong:} Opt-out of personalization is a DSA provision
\item \textbf{Why C is wrong:} The GDPR concerns personal data and the AI Act concerns AI risk
\item \textbf{Why D is wrong:} Self-preferencing is a DMA matter
\end{itemize}

\textcolor{myred}{\textbf{EXAM TRAP:}} DMA = gatekeepers and competition; DSA = platform content and transparency!
\end{frame}

\begin{frame}[shrink=6]{Q50: US vs China Approaches - Advanced}
\fontsize{6.5}{7.5}\selectfont
\textbf{QUESTION: Which comparison of AI governance in the US and China is accurate?}

\vspace{0.1cm}
\textbf{OPTIONS:}
\begin{enumerate}[A.]
\item The US has a comprehensive federal AI law; China relies on voluntary guidance
\item The US has no federal AI law but uses an executive order and the voluntary NIST AI RMF (govern, map, measure, manage); China has binding rules such as the 2023 Generative AI Measures, PIPL, and algorithm recommendation provisions
\item Both rely solely on voluntary frameworks
\item NIST AI RMF is mandatory certification
\end{enumerate}

\vspace{0.1cm}
\textcolor{myblue}{\textbf{ANSWER: B}}

\vspace{0.1cm}
\textbf{DETAILED EXPLANATION:}
\begin{itemize}
\item \textbf{Why B is correct:} The RMF is flexible and non-prescriptive. China's measures impose duties on generative AI providers and mandate data minimization, consent, and transparency.
\item \textbf{Why A is wrong:} The module says the US has no federal AI law
\item \textbf{Why C is wrong:} China has enacted binding rules
\item \textbf{Why D is wrong:} The RMF is voluntary
\end{itemize}

\textcolor{myred}{\textbf{EXAM TRAP:}} US: EO plus voluntary RMF; China: binding rules!
\end{frame}

\end{document}
